% ECONOMICS_PROBLEM: {"id": "D91-578", "title": "Empirically disciplined bounded rationality", "jel_code": "D91", "source_id": "OP-0003"}
% !TeX program = xelatex
\documentclass[11pt,a4paper]{article}
\usepackage[margin=23mm]{geometry}
\usepackage{fontspec}
\setmainfont{Latin Modern Roman}
\usepackage{amsmath,amssymb,mathtools}
\usepackage{xcolor,enumitem,booktabs,longtable,array,xurl,hyperref}
\definecolor{muted}{HTML}{53636C}
\hypersetup{colorlinks=true,urlcolor=blue,linkcolor=blue}
\setlength{\parindent}{0pt}
\setlength{\parskip}{5pt}
\newcommand{\R}{\mathbb R}
\newcommand{\E}{\mathbb E}
\newcommand{\N}{\mathbb N}
\newcommand{\ind}{\mathbf 1}
\newcommand{\Var}{\operatorname{Var}}
\newcommand{\Cov}{\operatorname{Cov}}
\newcommand{\supp}{\operatorname{supp}}
\DeclareMathOperator*{\argmin}{arg\,min}
\DeclareMathOperator*{\argmax}{arg\,max}
\newcommand{\entrymeta}[3]{{\sffamily\small\color{muted}\textbf{\texttt{#1}}\quad|\quad #2\par #3\par}}
\newcommand{\Problemref}[1]{\texttt{#1}}
\newcommand{\JELref}[2]{\texttt{#2} (JEL \texttt{#1})}
\begin{document}
\section*{Empirically disciplined bounded rationality}
\textbf{Problem ID:} \texttt{D91-578}\quad \textbf{JEL:} \texttt{D91}\par
% BEGIN ECONOMICS STATEMENT
\label{problem:OP-0003}
\label{atlas:prob:M3}

\noindent\textbf{Original atlas ID:} M3. \textbf{Formulation:} editorial specification of a research family.

\subsection*{Definitions and assumptions}

Let $\mathcal G$ be a declared population of finite games, with utilities normalized by a fixed experimental convention and observable covariates $x$. A structural model $m\in\mathcal M$ maps $(G,x,\alpha,\theta_m)$ into a probability distribution $p_m(\cdot\mid G,x,\alpha,\theta_m)$ on actions, where $\alpha$ is a subject characteristic with population law $F_m$ and $\theta_m$ is shared across games. For example, logit quantal response imposes
\[
\sigma_i(a)=\frac{\exp\{\lambda_i\,\mathbb E_{\sigma_{-i}}u_i(a,a_{-i})\}}
 {\sum_{b\in A_i}\exp\{\lambda_i\,\mathbb E_{\sigma_{-i}}u_i(b,a_{-i})\}},\quad\lambda_i\ge0.
\]
This is a fixed-point condition and may require an explicit equilibrium selector. Cognitive-hierarchy models instead specify a distribution of reasoning depths and beliefs about lower-depth behavior.

\subsection*{Formal research question}

Determine whether a finite-dimensional model family admits identified shared parameters and improves predictive performance across games drawn from a prespecified target distribution $Q$. Define the out-of-game log-score risk
\[
\mathcal R_Q(m,\theta_m,F_m)
 =\mathbb E_{(G,x)\sim Q}\mathbb E_{a\sim P_{G,x}}
 \left[-\log\int p_m(a\mid G,x,\alpha,\theta_m)\,F_m(\mathrm d\alpha)\right],
\]
where $P_{G,x}$ is the population action law and candidate probabilities must be positive on its support. Characterize observational equivalence, estimate risk using games withheld from fitting, and test restrictions that prevent game-specific parameters from absorbing every discrepancy. Report the identified set when the shared parameters or depth distribution are not uniquely identified.

\subsection*{Interpretation and scope}

The target is transportable prediction, not a theorem that one model describes all humans in all games. Utility normalization matters because payoff rescaling changes the effective logit precision unless its parameter is transformed accordingly. Subject heterogeneity, repeated-game experience, and equilibrium selection within quantal response must be represented in the likelihood rather than silently omitted. A within-subject game battery can distinguish stable subject effects from population compositional changes; assignment and attrition must still be handled. Comparing models by training likelihood alone does not answer the question. The held-out game distribution, score, benchmark, and permitted parameter dimension should be declared before collecting confirmatory outcomes.

\subsection*{Source audit and status}

McKelvey--Palfrey (1995), Camerer--Ho--Chong (2004), and Crawford--Costa-Gomes--Iriberri (2013) are verified with their exact titles and journal records below. The first develops quantal-response equilibrium, the second cognitive hierarchy, and the third reviews nonequilibrium strategic thinking. Existing fits and comparisons are substantive partial answers; none establishes universal parameter invariance. This entry is therefore an empirical modeling agenda with an editorial validation and identification target. The absence of a universal cross-game model in the cited papers is not a certified mathematical impossibility or proof that the specified empirical comparison remains unperformed in every relevant population.

\subsection*{References}

\begin{enumerate}

\item[S1] Richard D. McKelvey and Thomas R. Palfrey (1995). \emph{Quantal Response Equilibria for Normal Form Games}. Games and Economic Behavior 10(1), 6--38. \url{https://authors.library.caltech.edu/records/4fh4n-4c340}. \textit{Locator:} abstract; fixed-point definition and estimation. \textit{Establishes:} Existence and empirical estimation of quantal-response equilibria. \textit{Does not establish:} context-invariant cognitive parameters.

\item[S2] Colin F. Camerer, Teck-Hua Ho, and Juin-Kuan Chong (2004). \emph{A Cognitive Hierarchy Model of Games}. Quarterly Journal of Economics 119(3), 861--898. \url{https://academic.oup.com/qje/article-abstract/119/3/861/1938841}. \textit{Locator:} abstract; model and experimental applications. \textit{Establishes:} A parsimonious nonequilibrium reasoning model. \textit{Does not establish:} uniform prediction across all strategic environments.

\item[S3] Vincent P. Crawford, Miguel A. Costa-Gomes, and Nagore Iriberri (2013). \emph{Structural Models of Nonequilibrium Strategic Thinking: Theory, Evidence, and Applications}. Journal of Economic Literature 51(1), 5--62. \url{https://www.aeaweb.org/articles?id=10.1257/jel.51.1.5}. \textit{Locator:} abstract and review of level-k evidence. \textit{Establishes:} Evidence and applications of structural nonequilibrium models. \textit{Does not establish:} a theorem certifying the atlas prediction problem as open.

\end{enumerate}

\subsection*{Common conventions: Source mathematical conventions}

Unless an entry states otherwise, agent and firm sets are finite. State and action spaces are measurable, strategies and outcome maps are measurable, probability laws are specified, and expectations exist for the targets under discussion. Infinite-horizon discount factors lie in $(0,1)$ and discounted objectives are finite. Norms, tolerances and complexity input sizes must be specified for computational comparisons. Constants or primitives that appear locally are inputs, not empirical facts asserted by this catalogue.

\textbf{Identification.} A statistical model $m\in\mathcal M$ implies an observable law $P_m$ and target $\theta(m)$. At law $P$, its identified set is
\[
\Theta(P;\mathcal M)=\{\theta(m):m\in\mathcal M,\ P_m=P\}.
\]
Point identification holds when this set is a singleton, or equivalently when $P_{m_1}=P_{m_2}$ implies $\theta(m_1)=\theta(m_2)$ for all compatible models. Sharp bounds mean bounds attained, or approached when the boundary is not attained, by compatible models. Writing this set is a definition, not a solution: the substantive task is to establish informative restrictions and derive or estimate the set. An unrestricted-confounding model may give only trivial bounds.

\textbf{Inference.} Identification, estimability and valid uncertainty quantification are separate questions. Fix $\alpha\in(0,1)$. A confidence procedure $C_n$ over a declared law class $\mathcal P$ has asymptotic uniform coverage at level $1-\alpha$ when
\[
\liminf_{n\to\infty}\inf_{P\in\mathcal P}\Pr_P\{\theta(P)\in C_n\}\geq1-\alpha.
\]
For set-identified targets the coverage event must be defined separately, for example coverage of the full identified set. If an entry uses identification-indexed models instead of uniquely indexed laws, the same coverage convention is applied over those models. Pointwise limits do not establish uniform validity, and asymptotic validity does not guarantee finite-sample precision. Sample sizes, dependence structures and weak-identification sequences are part of each application.

\textbf{Counterfactuals.} $Y_i(a)$ is a potential outcome under a specified intervention, including its timing, implementation and versions. It is not $Y_i$ conditional on voluntarily choosing $a$. A full intervention vector permits interference; shorthand scalar treatment notation does not silently impose no interference. Assignment, selection, attrition, anticipation and measurement must be specified. A claim about an instrument's local effect is not automatically a population policy effect.

\textbf{Economic equilibrium.} A counterfactual model includes best responses, constraints, feasibility, market clearing, state transitions and expectations consistent with the stated information. When several equilibria exist, either a declared selector is an input or the target is set-valued. Comparisons across policy regimes must use the stated selection convention. Reduced-form relationships are not presumed invariant after an equilibrium-changing intervention.

\textbf{Optimization and welfare.} Optimization is over the stated feasible set. If attainment is not established, $\sup$ or $\inf$ and $\varepsilon$-optimal choices replace an asserted maximizer or minimizer. For example, with value $V=\sup_{a\in\mathcal A}W(a)<\infty$, an $\varepsilon$-optimal set is $\{a\in\mathcal A:W(a)\geq V-\varepsilon\}$ for $\varepsilon>0$. Benefits, costs, risk and health or environmental outcomes can be aggregated only using declared units and conversion weights. Interpersonal utility weights, the treatment of fiscal transfers and foreign profits, discounting, and the behavioral welfare criterion are explicit inputs. A comparative-static derivative additionally requires existence and appropriate differentiability of the selected counterfactual.

\textbf{Sources.} Local labels [S1], [S2], and so forth restart in every question. Locators identify the relevant abstract, model, section or result at the level actually checked; a bibliographic or abstract-level verification is not represented as inspection of an entire proof. Working papers are distinguished from later publications and changing versions. Source summaries are paraphrased. All URL checks and editorial formulations refer to the audit date stated above.
% END ECONOMICS STATEMENT
\end{document}
